\hypertarget{conclusions}{%
\chapter{Conclusions}\label{conclusions}}

The question was why the transfer of these assets from international to
indigenous operators produced such different results, and what should
follow from it. The answer is that capability-related factors best
account for the differences, which do not track the capital each
acquirer raised, the fiscal terms it faced, or the regulatory regime it
worked under. The gap that produces this divergence is built into the
structure of the transfer rather than into any single operator. What the
evidence warrants is a mandatory, asset-specific assurance mechanism for
high-risk producing-asset transfers.

Behind that answer is a particular way of seeing the asset. A producing
field is a sociotechnical system, the physical facilities and the
operating organization that runs them working as one, and a divestment
that conveys the facilities while leaving the organization behind
separates the system from the capability that kept it safe and
productive. In those terms, the cases reveal a control failure in how
the transfer is built, one whose cost falls largely outside the sale
(Figure 2.1). The framework developed in Chapter 5 is a control
structure designed to close that loop.

The answers to the research questions come first (Section 6.1), followed
by why the gap calls for a mandatory rather than a voluntary response
(Section 6.2). The recommendations that follow from it are set out next
(Section 6.3), then the limits of the findings (Section 6.4) and the
boundaries of their generalizability (Section 6.5). The contribution is
stated last (Section 6.6).

\hypertarget{answering-the-research-questions}{%
\section{Answering the Research
Questions}\label{answering-the-research-questions}}

This thesis opened with three questions, and the case analysis answers
each (Chapter 4). First, the differentiated outcomes are most consistent
with capability-related factors: the acquirers separated according to
the operating capability each inherited, brought, contracted, or
embedded before performance depended on it. Second, the evidence
contradicts fiscal regime effects for these cases and rejects capital
constraints as sufficient for the funded operators, finds regulatory
inadequacy and political economy each enabling but insufficient on its
own. It leaves capability-related factors as the explanation most
consistent with the integrity outcomes. Third, because the explanation
is capability-related, the remedy is a three-actor Capability Assurance
Framework that makes operating capability observable, testable, and
enforceable at the point of transfer, with the seller disclosing
requirements, the buyer demonstrating absorption, and the regulator
verifying both.

Across the two outcome dimensions, production performance and
operational integrity, the acquirers did not converge: they separated
into distinct trajectory groups, and the separation lines up most
closely with the operating capability each acquirer inherited, brought,
contracted, or embedded before performance depended on it. The
separation does not track the capital each raised, the fiscal terms each
faced, or the regulatory regime each worked under.

Capital constraints, the common explanation, are rejected as sufficient.
On the OML 29 critical experiment, the operator that suffered the most
serious integrity failure had funded and delivered production growth
beforehand, and an operator able to raise and deploy capital for
production cannot have its later integrity failure explained by capital
scarcity alone. A separate body of evidence points to capability: no
standing well-control contract or matching contingency plan when the
Santa Barbara well blew out, specialist intervention engaged only after
the fact, and a suspended, undecommissioned well the operator was
unready to control. The pattern points to capability, with the national
spill agency recording that the event exceeded its internal capacity
(Section 4.3). The two readings rest on different evidence.

Regulatory inadequacy is an enabling condition, not the primary driver.
The test is differentiation: were a weak regime the cause, acquirers
under it would perform alike, yet under the same pre-2024 regime that
governed Aiteo, Seplat and the HEOSL operatorship sustained production
without an equivalent integrity failure (Section 4.5). A regime common
to all cannot explain outcomes that varied sharply among them; what it
does is let the consequences of a capability shortfall go uncaught: it
is the setting in which a capability shortfall produces harm, while the
shortfall itself originates in the transfer.

Fiscal regime effects are rejected on timing: every primary transaction,
and the capability gaps that produced the failures, predate the
Petroleum Industry Act of August 2021, and the divergence under
identical pre-Act terms reinforces the point; whether the Act shapes
future transactions is left open. Political economy explains which firms
acquired and kept operatorship, and the production gains that
connections secure by suppressing theft, but it does not reach the
case-level capability to sustain integrity or contain a well-control
crisis. Why two firms on comparable terms then diverged in whether they
could run their assets safely is a question it leaves open.

The within-asset comparison on OML 30 is the strongest evidence against
the obvious competing reading, that outcomes track asset quality.
Operatorship passed from NPDC to HEOSL on the same wells, reservoirs,
and facilities, and performance improved under the new operator, which
is hard to reconcile with an asset-quality account. The population-level
evidence, and the marginal-field program of 2003 to 2020 in which a
majority of licensed operators also failed to sustain their assets, are
consistent corroborating context that does not, by itself, carry the
verdicts.

The five explanations resolve as follows. Capital constraints are
rejected as sufficient; fiscal regime effects are rejected for the cases
examined, the transactions and capability gaps predating the PIA.
Regulatory inadequacy is an enabling condition, insufficient on its own.
Political economy accounts for commercial survival and continued
participation but is insufficient on its own for operational integrity.
Capability-related factors are most consistent with the combined
pattern, with the OML 30 comparison strengthening that reading against
an asset-quality explanation. That answers the first two research
questions and grounds the design argument that follows.

\hypertarget{the-structural-case-for-mandatory-intervention}{%
\section{The Structural Case for Mandatory
Intervention}\label{the-structural-case-for-mandatory-intervention}}

The differentiated outcomes are not reducible to individual
mismanagement, to be corrected one operator at a time. They follow from
how the transaction itself is built. At the point of divestment the
information about an asset's embedded operating capability sits with the
seller, the capacity to absorb that capability sits with the buyer, and
the full social cost of a capability shortfall falls on neither. The gap
is structural, a control failure in the design of the transfer, not a
defect peculiar to any one acquirer.

Three features of the transfer make it so, each set out in the economics
of the transaction (Section 2.5). Information asymmetry comes first. The
seller knows what its operating organization does to hold an aging asset
within safe bounds; a buyer cannot read that from a data room, and the
regulatory screen at consent does not reach the asset-specific
capability at issue. The party that knows the most about the embedded
capability has little incentive to surface it, since fuller disclosure
can lower the price, slow the close, or create transition obligations,
and the party that carries the operational risk cannot verify what it is
buying. Absorptive capacity is the second, and it sits on the buyer's
side. A seller willing to hand over everything it knows still cannot
make that knowledge land in an organization without the prior expertise
to receive and embed it. The capability has to be present before
performance depends on it, which is the line the cases divided along.
The third is the externality. When an operator cannot contain a
well-control event or sustain integrity, the cost falls on host
communities and on the Nigerian state, parties outside the sale and
absent from its terms.

These three features have one consequence in common. Voluntary transfer
cannot solve a problem whose critical information is held by the seller,
whose absorption depends on the buyer, and whose failure costs are borne
by communities and the state. Voluntary instruments do not escape it.
The industry's operating and process-safety standards and the
civil-society National Principles for Responsible Petroleum Industry
Divestment describe the right content, but they attach no consequence to
ignoring it and leave the underlying incentives untouched (Section
5.5.3). Each describes what a responsible transfer should look like.
None compels it.

The instruments that do carry legal force have a different weakness: the
ones meant to protect the exposed parties are funded out of the
operations that fail. The Host Communities Development Trust is funded
by a contribution equal to three percent of the operator's actual annual
operating expenditure (PIA Section 240). When an operator falters and
operations contract, that contribution contracts with them, leaving the
fund thinnest when the hazard is greatest. The decommissioning
obligation runs the same way from the other direction. Where the
Decommissioning and Abandonment Fund is left to accrue across the
producing life of the asset rather than secured at transfer (Sections
232 and 233), an operator that fails early has set aside only a fraction
of what closing its wells will cost. Both mechanisms are necessary, and
each delivers least when it is needed most.

The structure points to a single remedy, security taken at the front of
the transaction rather than left to accrue. The Act already provides an
analogous structure: its environmental remediation provision requires a
financial contribution the Commission can reach (Section 103), fixed at
the point of operation rather than dependent on the operator's later
solvency. Extending that logic to operating capability and to community
exposure is the work of the recommendation developed below. The deeper
reason the transaction cannot correct itself is that the parties who
bear the residual cost are not at the table. Host communities and the
state cannot bargain ex ante over terms they never see, so the
externality cannot be priced into a private deal and must be placed
there by the body that licenses the transfer.

The third research question asked what institutional design follows from
these findings, and how responsibility should be allocated across the
seller, the buyer, and the regulator. A voluntary framework would
suffice only if market forces or the existing rules already secured
capability transfer. They do not. Outcomes diverged sharply under a
common regulatory regime, the transacting parties hold neither the
information nor the incentive to close the gap on their own, and the
full social cost of failure falls outside the sale. The rules already in
force do not reach the asset-specific operating-capability question at
consent. The Commission's 2024 reforms created a technical-capacity
route, but it is assessed at the corporate level, and a corporate screen
can pass an operator that has not demonstrated the asset-specific crisis
readiness the cases show to be decisive. The provisions that reach past
consent are triggered by failure rather than preventing it. What the
evidence warrants is a mandatory, asset-specific assurance mechanism for
high-risk producing-asset transfers, the control structure designed to
close the severed loop (Figures 5.1 and 5.2), allocating disclosure,
demonstrated absorption, and verification across the seller, the buyer,
and the regulator.

\hypertarget{prioritized-recommendations}{%
\section{Prioritized
Recommendations}\label{prioritized-recommendations}}

The full framework, with its three roles, five lifecycle phases, and
risk tiers, is developed in Chapter 5. For implementation, six
recommendations follow from it, ordered from before consent to after
close, each re-specifying an instrument the Nigerian regime already
contains rather than adding a new one. They carry the framework's
disclosure floor and its five structural levers (Section 5.6.3) into
instrument form. Together they put the framework's three roles into
operation: what the seller must disclose, what the buyer must
demonstrate, and what the regulator must verify and enforce. They are
risk-tiered: a disclosure, data, and continuity floor applies to every
transfer, while staged operatorship, technical-operator coverage, and
intensified post-transfer verification apply to the highest-risk assets
and lowest-capability buyers. Disclosure and continuity come first,
because the recommendations that follow depend on them.

\begin{enumerate}
\def\labelenumi{\arabic{enumi}.}
\item
  A regulator approving the transfer of a high-risk producing asset
  should require, as a condition of consent, that the seller disclose
  the asset's condition and the operating requirements it imposes, and
  that the buyer demonstrate the operating capability and
  crisis-response readiness to meet them. This follows from the
  documented condition that the information determining whether an aging
  asset can be held within safe bounds sits with the seller and cannot
  be read from a data room. In Nigeria the requirement attaches to the
  ministerial consent already mandatory for assignment of an upstream
  license or lease (PIA Section 95) and to the due-diligence requirement
  under the Nigeria Upstream Petroleum (Assignment of Interest)
  Regulations 2024 (Reg. 7). It makes asset-specific disclosure and a
  demonstrated well-control and crisis-response capability explicit
  conditions of consent. Neither is left to the buyer's discretion after
  close. Aiteo's transfer closed with no documented well-control
  contingency for OML 29 and no standing well-control contract, the gap
  that surfaced in the Santa Barbara blowout (Section 4.3). A
  disclosure-and-readiness condition at consent is the point at which
  that gap is still visible and correctable.
\item
  The transfer should carry binding terms that keep an aging asset's
  operating knowledge available to the buyer through and beyond the
  handover. That knowledge holds the asset within safe bounds, and it
  resides in the people who hold it, the procedures that codify it, and
  the working practices through which it is applied. Hardware transfers
  at completion; the embedded operating knowledge does not, and tacit
  knowledge in particular leaves with its holders. Nigeria already
  requires a succession plan and a multi-year understudy of incumbent
  expertise under the Nigerian Oil and Gas Industry Content Development
  Act (NOGICD, Section 31). That mechanism was built for
  expatriate-to-Nigerian succession. Re-specified toward seller-to-buyer
  transfer at divestment, and written into the assignment agreement
  under Section 95 as binding terms for retention, secondment, data, and
  understudy rather than negotiated ones, it would function as an
  obligation to carry the operating knowledge across the change of
  hands. The HEOSL recovery on OML 30, where a technical operator held
  that knowledge through the transition, shows what such terms are meant
  to secure.
\item
  Consent should depend on the buyer demonstrating the organizational
  capacity to receive and embed the disclosed capability, not on
  financial and legal standing alone. Disclosed knowledge cannot land in
  an organization that lacks the prior expertise to absorb it, which is
  the line the cases divided along. This re-specifies the
  technical-capacity criterion already within the consent assessment
  (Section 95; Reg. 7), raising it from a corporate-level screen to an
  asset-specific demonstration. For Tier 1 assets, the highest-risk
  tier, the demonstration should normally include a qualified technical
  operator with a record on comparable operations, drawing on the
  technology-transfer and partnering route the content law already
  contemplates (NOGICD Sections 43-47). In Aiteo, the
  leadership-background evidence pointed to downstream trading rather
  than comparable upstream operating depth. Seplat brought that depth,
  through its founding team and its technical partnership with Maurel et
  Prom. That contrast is the case basis for testing absorptive capacity
  before consent, not discovering its absence after.
\item
  For high-risk transfers, operatorship should pass in stages rather
  than at a single completion date, with critical operating functions
  passing only after the buyer has demonstrated readiness for them.
  Time-bounded transition support, technical-operator coverage, or
  secondment holds the asset within safe bounds during the transition,
  without leaving the seller indefinitely responsible for an asset it
  has sold. An abrupt handover concentrates the capability discontinuity
  at the moment of greatest exposure, before the buyer has shown it can
  run the asset. In Nigeria staged operatorship is a condition the
  Commission can attach to consent under Section 95, supervised under
  its standing authority to monitor and enforce the terms of leases and
  licenses (Section 7(d)). It requires no new instrument. The phased
  technical-operator arrangement through which OML 30 was returned to
  stable operation is the design pattern. The single-date handover at
  Aiteo, into an organization without demonstrated upstream operating
  depth, is the failure that staged operatorship is built to prevent.
\item
  The regulator should verify operating capability on a continuing basis
  after transfer, prompted by rule-based triggers, production or
  integrity deviations, missed maintenance intervals, the loss of key
  operating personnel, rather than by discretion alone. A capability
  shortfall invisible at consent becomes visible only in operation, and
  an unmonitored latent exposure is the most dangerous kind. Nigeria's
  regulator already holds the powers this requires: to monitor and
  enforce compliance with the terms of leases and licenses (Section
  7(d)), and to carry out audits and investigations of licensees
  (Section 7(g)), with administrative penalties available for
  contravention (Section 231). Binding those powers to defined triggers
  connects an integrity finding to a graduated, pre-defined consequence
  rather than to discretion, subject to the operator's opportunity to
  respond. The powers already exist; the triggers are what is missing.
  The well left suspended and undecommissioned on OML 29 from 2015,
  unmonitored until it failed in 2021 (Section 4.3), is the exposure
  continuing verification is built to catch.
\item
  The financial security for decommissioning and for community exposure
  should be taken at the point of transfer rather than allowed to accrue
  out of the operations that may fail. An instrument funded from ongoing
  operations delivers least when the operator falters, which is the
  moment it is most needed. The remedy is to secure it at consent. The
  amount should scale to the asset's risk and inherited liabilities.
  This removes the dependency the existing instruments carry, the Host
  Communities Development Trust's three-percent-of-operating-expenditure
  base (Section 240) and the decommissioning fund's accrual across the
  producing life of the asset (Sections 232 and 233). It extends the
  upfront structure the Act already uses for the environmental
  remediation contribution (Section 103) to decommissioning security and
  community-readiness provisioning fixed at consent. The
  undecommissioned well on OML 29 and the 41 communities exposed by the
  blowout are the costs this provisioning is meant to cover before, not
  after, an operator's capacity to pay has eroded.
\end{enumerate}

The six divide by where they attach. Recommendations 1 through 4 attach
to the consent gate, the one-time approval at transfer. They therefore
apply only to transfers not yet completed, the divestments still to come
in a wave of IOC onshore exits that is not finished, and cannot be
re-run on a transaction already closed. The consent-stage screening they
require is not hypothetical: the Commission applied a version of it when
it initially withheld consent for the Shell-to-Renaissance transfer in
2024 on the ground that the consortium had not demonstrated the
requisite technical capacity. The Commission granted consent later that
year (The Energy Year, 2025; BusinessDay, 2025). Recommendations 5 and 6
attach to powers and obligations that bind every current operator
continuously: the Commission's standing supervisory powers (Sections
7(d) and 7(g)) and the standing fund obligations under Sections 232-233
and 103. They therefore reach the assets already transferred, Aiteo and
Eroton included. Applying them to current operators enforces standing
obligations rather than reopening closed consent decisions. It would not
re-run the original transfer; it would apply continuing powers and
standing obligations to operating risks that are still live. That is
where the exposure now sits.

None of this builds a second regulatory track alongside the one the
Commission introduced in 2024. That route assesses capacity at the
corporate level; the recommendations re-specify the same consent gate
and the same standing powers to the asset-specific operating capability
and crisis readiness the cases show to be decisive. Nor do they depend
on settling whether a given operator's capability came from what it
inherited at the transfer or from its own organizational quality. Each
tests for demonstrated operating capability directly, whatever its
source, a limit taken up in the next section (Section 6.4). These
measures work as a set, not a menu. For a new transfer, requiring
disclosure at consent without the matching post-transfer verification
would let a shortfall that clears the gate go uncaught, so the
consent-stage and post-transfer measures belong together rather than
sequenced years apart. What the six do, operating together, is make the
operating capability an asset requires observable before consent,
testable at the gate, and enforceable after transfer. They are warranted
by the conditions the cases document, not by evidence of their effect,
which no implementation has yet produced; the claim is that they make
capability assurable, not that they will prevent the next incident or
raise national production.

\hypertarget{limitations}{%
\section{Limitations}\label{limitations}}

The findings carry weight within stated bounds, and naming those bounds
is part of stating the findings accurately. Each limitation below
constrains the scope of a particular claim, and none overturns the
central finding.

The deepest of these concerns the regulator. Verification rests with the
same Commission whose limited reach was an enabling condition of the
failures (Section 6.1), which invites the objection that an institution
unable to catch the old shortfall cannot be relied on to operate the new
safeguard. The design depends less on regulatory judgment than the
status quo does. It re-specifies the technical-capacity screen from
corporate-level standing to asset-specific readiness, shifts the burden
of proof to the buyer to demonstrate capability, and ties consequences
to rule-based triggers rather than to discretion. The parts most exposed
to weak capacity, site visits and post-transfer audits, can be
risk-triggered rather than applied to every transfer, while disclosure
and transfer-agreement content can be enforced as documentary
requirements. The limit is real even so. If the Commission lacks the
capacity, independence, data access, or political backing to apply the
design, the framework cannot function as intended. The framework and the
building of regulatory capacity are jointly necessary, not substitutes,
and a phased adoption that begins with disclosure and transfer terms and
builds toward full verification is the realistic path in a constrained
environment.

A second limit concerns the source of the capability the cases turn on.
The evidence cannot fully separate whether a given operator's capability
came from what it inherited at the transfer or from organizational
quality it already held, so the causal attribution to the inheritance
itself is bounded. The prescription does not depend on settling that
attribution. Each recommendation tests for demonstrated operating
capability at the point of transfer, whatever its origin (Section 6.3),
so the confound limits what can be claimed about cause without weakening
what the framework requires.

A related reading attributes the divergence to the assets rather than to
the operators, setting Aiteo's failure on a large and difficult asset
against acquirers that performed on easier ones. The within-asset
comparison on OML 30 is the direct answer (Section 6.1): the operator
changed while the asset was held fixed, and the outcome changed with it.
Where that holds, the asset-quality explanation is substantially
weakened, and the operator-capability finding stands.

The framework has not been implemented, and no implementation data yet
show that it improves performance, prevents incidents, or raises
production. The claim made for it is therefore the bounded one: it is
warranted by the regularities documented across the cases and by the
design logic, not proven by results it has yet to produce. What the
absence of outcome data constrains is any claim about effect. It does
not weaken the warrant, which rests on harms that are documented, that
fall on parties outside the sale, and that are difficult to reverse once
they occur (Section 6.2). Action is warranted under those conditions
even though the framework remains untested as an intervention. What the
framework changes, independently of any later results, is that the
capability the cases tie to failure must be shown at the point of
transfer rather than assumed.

Verification reaches the capability that is observable, demonstrable,
and verifiable. The non-contractible tacit knowledge that no instrument
can compel a seller to surrender stays with the individuals and teams
who hold it. For that reason the recommendations secure the conditions
under which knowledge transfers, retention, secondment, and a documented
handover, rather than the knowledge itself. The boundary limits what
verification can guarantee. It does not remove the value of verifying
what is observable, since several of the capabilities the cases turned
on, crisis-readiness arrangements, technical-operator depth, and
continuity of operating knowledge, were demonstrable before close.

The evidence is a set of six cases built around a single critical
experiment with cross-case validation (Chapter 3). A design of this kind
can corroborate an explanation and show it to be plausible across the
cases examined; it cannot establish that no counter-example exists, and
it does not prove a universal law. The finding is therefore stated in
calibrated form: capability-related factors are the explanation most
consistent with the observed pattern, not a demonstrated universal
regularity. What these cases support beyond their own setting is taken
up next (Section 6.5).

The reach of the consent-stage recommendations into transactions already
closed is partial by design. Because they cannot be applied to a
divestment that has already completed (Section 6.3), their effect is
bounded to pending and future transfers, which limits what they can
change without bearing on the diagnosis they follow from. The
standing-power recommendations do reach the assets already transferred,
so the framework is not inert for the operators now holding them.

These limits together fix where the thesis's claims stop. They bound the
causal attribution, the outcome claim, the reach of verification, and
the generality of the result, but they do not unsettle the diagnostic
finding or the design warrant. What stands is the case that mandatory,
asset-specific assurance is warranted for high-risk producing-asset
transfers under the conditions documented here.

\hypertarget{boundaries-of-generalizability-and-future-research}{%
\section{Boundaries of Generalizability and Future
Research}\label{boundaries-of-generalizability-and-future-research}}

How far these findings reach is a question of analytic transferability,
not demonstrated generalizability. The framework was developed on
Nigerian divestment cases between 2010 and 2024. It is calibrated first
to Nigeria and, analytically, to the wider institutional context set out
in Objective 3: petroleum-producing developing countries in which a wave
of divestment meets constrained regulatory capacity and an externality
placing the cost of failure on host communities and the state. To other
settings that share those conditions the framework may transfer, subject
to local validation. The basis is analytic rather than statistical,
because the cases were chosen to test competing explanations rather than
to represent a population (Chapter 3).

The comparator regimes do not validate the framework, but they show that
its core move, requiring capability to be shown before and after a
transfer rather than merely declared, is established practice in mature
petroleum jurisdictions, not an invention of this thesis. The United
Kingdom requires the buyer and seller to file a capability pack
documenting the buyer's financial and technical capability before it
consents to a license assignment. It also treats the fitness of those
who control the licensee as a continuing criterion, backed by a standing
power to require a change of control or to revoke the license. Norway
prequalifies operators before entry and requires a competent in-country
organization. Alberta assesses both parties to a transfer and ties the
security a buyer must post to that assessment (Section 5.5). What
differs is the target. Those regimes run the logic inside strong
institutions, whereas the framework adapts it to a setting where
regulatory capacity is thinner and the consequences of failure fall on
parties outside the sale, the harder case for which it is designed.

Two features of the evidence bound how far the finding extends. Some of
the Nigerian setting is specific to it: the theft and sabotage that
confound production outcomes at the sector level (Chapter 4) limit what
the production evidence can show, though they touch neither the
existence of the externality nor the need to verify capability. The
cases are also weighted toward divestments originating with Shell, which
narrows the seller-independence of the finding. That weighting does not
unsettle the core result, because the outcomes diverged sharply among
buyers who all acquired from the same seller, so the divergence cannot
be an artifact of seller identity. What it leaves open is whether the
same pattern recurs across a broader base of sellers.

The research agenda follows from these limits, not from a general call
for more study. The first and most important open question follows from
the absence of implementation data (Section 6.4): the framework is
warranted by the documented conditions, not tested by results. The major
post-2024 divestments, the Shell-to-Renaissance transfer and the recent
ExxonMobil/MPNU and Eni/NAOC transfers to Seplat and Oando, are closing
under the 2024 reforms and will produce the outcome record the cases
here could not. As that record accumulates, it moves the design from
warranted toward tested. A comparison between transfers screened for
asset-specific capability and those screened only at the corporate level
would then test, refine, or qualify it. These transactions are future
tests, not present evidence; none has yet run long enough to establish
medium-term production, integrity, and governance outcomes.

Four narrower questions follow from limits already stated. The
organizational-quality confound (Section 6.4) would be settled by
operator-level data that separates capability inherited at transfer from
capability a buyer already held, the attribution the cases could only
bound. The seller-independence question would be answered by cases drawn
from a wider set of divesting parties than the Shell-weighted record
here. A third question concerns the boundary of non-contractible tacit
knowledge (Section 6.4): how far retention, secondment, and a documented
handover can carry what no instrument can compel directly. The last is
one of scope, whether the framework holds in other petroleum-producing
developing countries where divestment, constrained regulatory capacity,
and externalized failure recur.

Each of these is a specific open question left by a stated limit, not a
gap to be filled for its own sake. With the evidence already in hand,
they mark the path by which a framework warranted on the Nigerian record
could be tested, extended, and, where the conditions recur, carried
beyond it.

\hypertarget{contribution-and-conclusion}{%
\section{Contribution and
Conclusion}\label{contribution-and-conclusion}}

The primary contribution of this thesis is a systems-design one, and it
reduces to a single reframing. What fails to transfer in a divestment is
the operating capability embedded in the organization that ran the
asset: the wells, the facilities, and the licenses change hands intact,
while the capability that kept them safe and productive stays behind.
Closing the gap that opens is a problem of institutional design, met by
a control loop that holds that capability observable, demonstrable, and
verifiable across the transfer. An upstream divestment is then the
transfer of a sociotechnical system, and the failures observed across
the cases sit in the control structure of that system. The empirical
finding, that capability-related factors best account for the
differentiated outcomes, underwrites the reframing; the reframing and
the design that follows from it are the contribution.

The diagnosis is grounded in a tradition that treats a producing asset
as a coupling of a technical subsystem and the social subsystem that
operates it (Trist and Bamforth, 1951; Section 2.2). Safe and productive
operation is an emergent property of the two together rather than of the
equipment alone. Unless the operating capability is deliberately carried
across, a divestment divides that system along the interface between
them, conveying the wells, the facilities, and the data while leaving
the operating organization behind. What that organization holds, from
the operating workforce to the contractor architecture, is the operating
enterprise itself, which enterprise-architecting treats as an object of
design in its own right (Nightingale and Rhodes, 2015; Section 5.2).
Through that lens, the failure has a precise location. The party that
knows what is being lost has no incentive to transfer it, the corrective
feedback that would catch a shortfall before harm is severed, and much
of the cost falls on parties outside the sale (Figure 2.1). That is a
control failure in the systems-theoretic sense, the absence of the
control actions and feedback that hold a hazardous process within safe
bounds (Leveson, 2011; Section 5.4).

The design follows from the diagnosis. The Capability Assurance
Framework is a control structure for the transfer: three mutually
constraining roles, seller disclosure, buyer absorption, and regulator
verification, that reconstitute the operating system across the
transaction by closing the loop the divestment severs (Figures 5.1 and
5.2). Each role supplies what the next must act on, so the structure
holds only when all three operate together, and it installs the
corrective function the market does not supply. The interlock sits at
the level of capability. A divestment re-architects the operating
enterprise only in part: it transfers the physical infrastructure intact
while leaving the knowledge that runs it to be reconstituted by the
buyer or lost. Because that gap is not captured by ordinary transaction
documents, it can remain unverified until it surfaces in operation, the
Aiteo pattern (Section 4.3). The buyer's absorptive capacity (Cohen and
Levinthal, 1990) is, in these terms, the receiving enterprise's capacity
to reconstitute that knowledge through the transition, and the
framework's verification step makes the difference observable before
handover rather than after failure.

The systems-design contribution is supported by three others, and none
displaces it. The first is theoretical, combining the
sociotechnical-systems tradition, enterprise-architecting,
systems-theoretic control, and the absorptive-capacity and
knowledge-transfer literatures into one account of a problem usually
treated in financial, fiscal, or regulatory terms. The second is
methodological: a discriminating-case design, built around a critical
experiment with cross-case validation and the explicit adjudication of
competing explanations (Chapter 3), adaptable to other divestment and
asset-transfer settings where the same alternatives compete. The third
is practical: the recommendations are attached to specific instruments a
regulator already holds (Section 6.3), which makes the design deployable
rather than only diagnostic.

The scope of this contribution is bounded by what the evidence supports.
The reframing and the design are warranted by the documented
regularities and the design logic, not by implementation results the
framework has yet to produce. The claim that holds is narrower and,
within the cases examined, secure. A divestment leaves the control loop
open when capability is not assured across it, and a framework of this
kind makes that capability observable, testable, and enforceable at the
point of transfer.

Yet the stakes reach beyond the transaction itself. A divestment is not
only a private sale; it is the transfer of responsibility for a public
hazard. When the control loop is left open, the cost of a capability
that does not survive falls not on the seller that exits or the buyer
that underperforms, but on the host communities and the Nigerian state
that never sat at the table. Closing the loop places the burden of
demonstrating capability on the parties to the transfer, where the cost
of failure says it belongs. Divestment, seen as a systems-design
problem, comes down to a single requirement: that the transfer of the
hardware not become the transfer of a hazard no one has agreed to carry.

\addtocontents{toc}{\protect\appendixtocsetup}
